\section{Free massive fields}\label{sec:free}

Here $d\ge1$, $n\ge2$, $L\ge n+1$, $\mathbb T=(\Z/L\Z)^d$, $\mu>0$ is a lattice
mass, $-\Delta$ is the lattice Laplacian and $Q=-\Delta+\mu^2$. For a positive
definite $Q$ on $\R^N$ let $\gamma_Q$ be the centred Gaussian law with
covariance $C=Q^{-1}$, and for $V\subset\{1,\dots,N\}$ let $\mathcal P_V$ be the
conditional expectation given the coordinates outside $V$. For the cubes
$V_x=x+\{0,\dots,n-1\}^d$, $x\in\mathbb T$, put
\begin{equation}\label{eq:ffvarpi}
\varpi(Q)=\inf_f\frac{n^{-d}\sum_x\|f-\mathcal P_{V_x}f\|^2}{\Var_{\gamma_Q}(f)},
\end{equation}
the infimum over nonconstant $f\in L^2(\gamma_Q)$, the best constant in
$\BP(n,\cdot)$. With $s=n\mu$ and $\mu=ma$, $n=\ell/a$, the parameter $s=m\ell$
is the mass times the physical size of the cube.

\begin{lemma}[linear functions suffice]\label{lem:chaos}
For a positive definite $Q$ and any family of coordinate sets $(V_x)$, the
infimum in~\eqref{eq:ffvarpi} equals its infimum over the linear functions
$\ell_h(\phi)=\ip h\phi$, and
\[
\|\ell_h-\mathcal P_V\ell_h\|^2=\ip{h_V}{Q_{VV}^{-1}h_V},\qquad
\Var(\ell_h)=\ip h{Q^{-1}h},
\]
with $h_V$ the restriction of $h$ to $V$ and $Q_{VV}$ the block of $Q$ on $V$.
\end{lemma}

\begin{proof}
The conditional law of $\phi_V$ given $\phi_{V^c}$ is Gaussian with covariance
$Q_{VV}^{-1}$, which gives the formula for $\ell_h$. Write $\phi=Q^{-1/2}\xi$ with
$\xi$ standard Gaussian. Conditioning on $\phi_{V^c}$ is conditioning on the
Gaussian subspace $K_V=\{\ip u\xi:u\in Q^{-1/2}\R^{V^c}\}$, and on the $m$-th
Wiener chaos, identified with the symmetric tensors $\mathrm{Sym}^m(\R^N)$, it
acts as $\pi_V^{\otimes m}$, with $\pi_V$ the orthogonal projection onto $Q^{-1/2}
\R^{V^c}$~\cite[Ch.~4]{Jan97}. The chaoses are orthogonal and invariant under
every $\mathcal P_V$, so it suffices to treat one chaos $m\ge1$. The projection
$I-\pi_V^{\otimes m}$ dominates $(I-\pi_V)\otimes I^{\otimes(m-1)}$, as both are
projections that commute and the range of the second lies in that of the
first. For a symmetric tensor $F$, therefore,
\[
\sum_x\|F-\pi_{V_x}^{\otimes m}F\|^2\ge\Bigl\langle F,\Bigl(\sum_x(I-\pi_{V_x})\Bigr)
\otimes I^{\otimes(m-1)}F\Bigr\rangle\ge\lambda\|F\|^2,
\]
where $\lambda$ is the least eigenvalue of $\sum_x(I-\pi_{V_x})$ on $\R^N$, which is
the infimum of the quotient over linear functions ($m=1$).
\end{proof}

\begin{theorem}[the exact constant]\label{thm:free}
Let $V=\{0,\dots,n-1\}^d$, $-\Delta_D$ its Dirichlet Laplacian and $\mathbf 1$ its
constant function. Then
\[
\varpi(-\Delta+\mu^2)=\frac{\mu^2}{n^d}\ip{\mathbf 1}{(-\Delta_D+\mu^2)^{-1}\mathbf 1},
\]
attained only at multiples of $\ell_{\mathbf 1}$ among linear functions, and
\[
\Bigl(\frac56\Bigr)^d\frac{s^2}{s^2+10d}\le\varpi(-\Delta+\mu^2)\le\frac{s^2}
{s^2+4dn^2\sin^2\frac\pi{2n+2}}\le\frac{s^2}{s^2+16d/9}.
\]
\end{theorem}

\begin{proof}
Since $L\ge n+1$, the block $Q_{V_xV_x}$ is the translate of $Q_{VV}=-\Delta_D+
\mu^2$; put $G_V=Q_{VV}^{-1}$. By Lemma~\ref{lem:chaos}, $n^d\varpi$ is the least
value of $\ip h{Sh}/\ip h{Q^{-1}h}$ with $S=\sum_xR_{V_x}^{\mathsf T}Q_{V_xV_x}^{-1}
R_{V_x}$, that is, the least eigenvalue of $Q^{1/2}SQ^{1/2}$. Translations commute
with $Q$ and $S$, so the plane waves $e_k(y)=e^{ik\cdot y}$ diagonalise it with
eigenvalues $\lambda_k=\hat q(k)\ip{e_k^V}{G_Ve_k^V}$, where $\hat q(k)=\mu^2+
\sum_i(2-2\cos k_i)$ and $e_k^V$ is the restriction to $V$.

Let $B$ map $f$ to the function $u\mapsto\sum_{w\sim u,\,w\notin V}f(w)$ on $V$.
Then $Q_{VV}R_Vf=R_VQf+Bf$, so $Q_{VV}e_k^V=\hat q(k)e_k^V+Be_k$, and pairing
$G_V$ of this with $e_k^V$ gives
\[
\lambda_k=n^d-\sum_{u,u'\in V}\sum_{w\sim u',\,w\notin V}G_V(u,u')\cos\bigl(k\cdot
(w-u)\bigr).
\]
Write $Q_{VV}=(2d+\mu^2)(I-K)$ with $K$ entrywise nonnegative and $\|K\|\le
2d/(2d+\mu^2)<1$; the Neumann series shows that $G_V$ has positive entries.
Hence $\lambda_k-\lambda_0=\sum G_V(u,u')(1-\cos(k\cdot(w-u)))\ge0$, with equality only
at $k=0$ (take $u=u'$ on a face and $w=u+e_i$). Since $Q\mathbf 1=\mu^2\mathbf 1$,
$\lambda_0=\mu^2\ip{\mathbf 1}{G_V\mathbf 1}$, which is the formula.

\emph{Upper bound.} $G_V\le(\lambda_1+\mu^2)^{-1}$ with $\lambda_1=4d\sin^2\frac\pi
{2n+2}$ the least eigenvalue of $-\Delta_D$, and $\|\mathbf 1\|^2=n^d$; then
$\sin t\ge2t/\pi$ gives $n^2\lambda_1\ge4dn^2/(n+1)^2\ge16d/9$.

\emph{Lower bound.} By the variational formula $\ip{\mathbf 1}{G_V\mathbf 1}\ge
\ip{\mathbf 1}w^2/\ip w{Q_{VV}w}$, take $w(u)=\prod_i\tau(u_i)$ with the parabolic
profile $\tau(s)=6(s+1)(n-s)/(n(n+1)(n+2))$ on $\{0,\dots,n-1\}$, zero
outside. Then $\sum\tau=1$, $S_2=\sum\tau^2=\alpha_n/n$ and $D_2=\sum_{s\in\Z}
(\tau(s+1)-\tau(s))^2=\beta_n/n^3$ with
\[
\alpha_n=\frac{6(n^2+2n+2)}{5(n+1)(n+2)}\le\frac65,\qquad
\beta_n=\frac{12n^2}{(n+1)(n+2)}=10\rho_n\alpha_n,\quad
\rho_n=\frac{n^2}{n^2+2n+2}<1,
\]
by the power sums $\sum_{u=1}^nu(n+1-u)=n(n+1)(n+2)/6$, $\sum_{u=1}^nu^2(n+1-u)^2=
n(n+1)(n+2)(n^2+2n+2)/30$ and $\sum_{u=0}^n(n-2u)^2=n(n+1)(n+2)/3$. The Dirichlet
form of the product is $\ip w{-\Delta_Dw}=dS_2^{d-1}D_2$ and $\ip{\mathbf 1}w=1$, so
\[
\varpi\ge\frac{\mu^2}{n^d(\mu^2S_2^d+dS_2^{d-1}D_2)}=\frac1{\alpha_n^d+d\alpha_n^{d-1}
\beta_ns^{-2}}=\alpha_n^{-d}\frac{s^2}{s^2+10d\rho_n}\ge\Bigl(\frac56\Bigr)^d
\frac{s^2}{s^2+10d}.\qedhere
\]
\end{proof}

At $n=2$ the two bounds coincide with $s^2/(s^2+4d)$. For $d=4$ the lower
bound is $0.0118$ at $s=1$ and $0.0438$ at $s=2$; the upper bound vanishes as
$s\to0$, so a massless free field satisfies no cube condition uniform in the
lattice spacing. The same proof, with the Hodge identity of the torus, gives
the cube condition for the free massive vector (Proca) field on gauge-invariant
functions, with the same lower bound and the upper bound $s^2/(1+s^2)$
\cite{BBLRS26}.

